\section{Review of RPL Objective Functions Load Balancing}
\subsection{Related Work}
Various objective functions have been proposed, which can be
categorised into 4 groups: single metric, composite and context aware.

OF0 and MRHOF fit into single metric category.

% TODO: maybe include more OF0 and MRHOF evaluation
\cite{pradeskaPerformanceAnalysisObjective2016} evaluate  the
performance of OF0 and MRHOF in COOJA Contiki
Additionally, ... evaluates the performance of the standardised OFs
and find performance of OF0 degrade upon increasing load

In \cite{kimQURPLQueueUtilization2015}, the authors argues that most
of the losses are due to queue loss. Thus they used queue
utilisation metric in objective function, called QU-RPL to deal with
dropped packets due to queue overflow. Additionally, they also
modified the trickle timer to reset when a node experiences many
consecutive queue loss. With the changes, they achive lower queue
loss ratio and a more balance tree. However, only queue metrics are
taken into consideration, without mentioning latency.

A new objective function: OF-FL is presented in
\cite{gaddourOFFLQoSawareFuzzy2014}, which employs fuzzy logic with
four metrics with the aim of supporting various application requirements.
The authors used four metrics hop count, end to end delay, battery
level and ETX for categorise link quality. Each metric follows fuzzy
logic to indicate link quality from awful to excellent. The
simulation results show that OF-FL achieve higher PDR while
conserving more energy. However, the process of encoding each metrics into quality label is done manually and subject to the biases of the doers. % TODO: more about limitation in this

The authors in \cite{abreuEnergyawareRoutingBiomedical2014} propose an energy-aware objective function in RPL for healthcare monitor application. The nature of this problem means that the sensors is mobile, thus affecting the link quality % TODO: should I include this paper

The authors in \cite{taghizadehCLRPLContextAwareLoad2018} proposed Context-Aware and Load Balancing RPL (CLRPL). They introduced a new metric called Context-Aware Routing Metric (CARM), which considers queue utilisation and the percentage remaining power of a entire parents' chain from root the the node instead of the immediate parent. CARM would cnosider low-power parent but high-power overall route instead of rejecting it like energy of a single parent.
% TODO: include limitation of this 


Weighted Random Forward RPL (WRF-RPL) was propsoed in \cite{acevedoWRFRPLWeightedRandom2021}, which would select parents randomly based on weighted score % TODO: more about this

In recent, there has been many attempts to employ machine learning
into OFs. Despite that, most of the new application of ML in RPL
focues on reinforcement learning
with little focus on supervised learning

In \cite{duenassantosQRPLQLearningBasedRouting2024}, the authorse
puts Q-Learning to OF, award the model when a packet is trasmitted
successfully and penalise when choosing a link that lead to packet loss.
% TODO: include more reinforcement learning algo
% TODO: include the metrics name

There are three researches on applying supervised ML to RPL OF from a
research group. In RPL+ , the authors utilise Decision Tree to
perform feature importance evaluation with
six metrics ETX, channel utilisation, density, through put, queue
utilisation and MAC losses on data gathered hop by hop from
simulation \cite{duenassantosRPLImprovedParent2022}. They set
the $Rank\_Increase$ fixed to 1, which can be regarded as hop count,
and use the feature
importance as coefficients in linear equation and select parent with
the lowest score. Even though they improve the performance of PDR and
delay, the random forest model is only for feature importance
analysis and the training data only consider the success chance of
forwarded by one hop instead of the whole path. Additionally, having
fixed $Rank\_Increase$ means that the OF would select parents with
the shortest hop count.

In \cite{mezherGNBRPLGaussianNaive2023}, the research group train a
Naive Bayes classifier to predict the probability of a packet is sent
successfully or not on the same training feature as RPL+. This assume
that the distribution of the features fall into Gaussian distribution
In this paper, they train on more data and integrate the model into
the OF. However the evaluation lacks consideration for OF0.

Subsequently, the group proposed ML-RPL, using CatBoost model to
learn from eight metrics namely hop count, ETX, MAC losses, density,
channel utilisation, throughput, queu utilisation and RSSI
\cite{santosMLRPLMachineLearningBased2023}. Similar to previous
researches, this model predicts the predict the probability of
successfull one hop transmission. The evaluation of the results lack
energy or cpu util evaluation % TODO: more about limiation


% TODO: more discussion on the limitation of these papers
